\subsection{Funzioni Test}
\label{subsec:TestRequests}

\subsubsection{run\_test}  
\label{subsubsec:run_test}

\paragraph{Descrizione}  
La funzione \textbf{run\_test} utilizza le domande del dataset per valutare il modello LLM selezionato, restituendo i risultati e le relative statistiche in formato strutturato.


\paragraph{Dettagli dell'endpoint}
\begin{itemize}
    \item \textbf{HTTP Method}: POST
    \item \textbf{Endpoint}: \texttt{/test}
    \item \textbf{Blueprint}: \texttt{api}
    \item \textbf{Response Format}: JSON
    \item \textbf{Nome funzione}: \texttt{run\_test}
    \item \textbf{Decoratori usati}: \texttt{@api.route}
\end{itemize}

\paragraph{Parametri di input}
\begin{itemize}
    \item \textbf{dataset$\backslash$questions}: la lista di domande da utilizzare per il test;
\end{itemize}

\paragraph{Struttura della risposta}
La risposta è un oggetto JSON composto da:
\begin{itemize}
    \item \textbf{generated\_answers}: le risposte generate dal modello;
    \item \textbf{scores}: punteggi assegnati in base al confronto con le risposte corrette;
    \item \textbf{metrics}: eventuali statistiche riassuntive (media, deviazione standard, ecc.).
\end{itemize}

\paragraph{Funzionamento}
\begin{itemize}
    \item Verifica la presenza dei dati nella richiesta e li valida;
    \item Esegue la funzione di valutazione passando le domande al modello LLM selezionato;
    \item Confronta le risposte generate con quelle corrette e calcola i punteggi;
    \item Restituisce un oggetto JSON contenente le risposte del modello, le valutazioni e i dati statistici.
\end{itemize}